\documentclass[UTF8,a4paper,12pt]{ctexart}

\usepackage[margin=2.35cm]{geometry}
\usepackage{amsmath,amssymb,bm}
\usepackage{booktabs,longtable,tabularx,array,multirow}
\usepackage{xcolor}
\usepackage{enumitem}
\usepackage{listings}
\usepackage{hyperref}
\usepackage{fancyhdr}
\usepackage{float}

\definecolor{reportgreen}{RGB}{36,92,68}
\definecolor{reportgray}{RGB}{245,247,246}
\hypersetup{
  hypertexnames=false,
  colorlinks=true,
  linkcolor=reportgreen,
  urlcolor=reportgreen,
  citecolor=reportgreen,
  pdfauthor={CUMCM 备赛小组},
  pdftitle={2020--2025 年 CUMCM A/B/C 题学习与论文写作报告}
}
\pagestyle{fancy}
\fancyhf{}
\fancyhead[L]{CUMCM A/B/C 题备赛学习报告}
\fancyhead[R]{2020--2025}
\fancyfoot[C]{\thepage}
\setlength{\headheight}{15pt}
\setlength{\parindent}{2em}
\setlength{\parskip}{0.35em}
\setlist{nosep,leftmargin=2em}
\lstset{
  basicstyle=\ttfamily\small,
  frame=single,
  breaklines=true,
  columns=fullflexible,
  showstringspaces=false
}

\title{\bfseries 2020--2025 年全国大学生数学建模竞赛\\A/B/C 题学习与论文写作报告}
\author{队长、组员1、组员2}
\date{2026 年 7 月 28 日}

\begin{document}
\maketitle

\begin{abstract}
本报告面向三人本科生队伍的全国大学生数学建模竞赛备赛。研究范围严格限定为
2020--2025 年 A、B、C 三类题，论文语料为工作区中 59 篇编号 PDF，共 2892 页；
D、E 题不进入统计。报告先按机理建模、综合优化和数据决策三条主线划分个人责任，
再把思维导图中的算法逐项落实为原理、步骤、代码接口、适用边界和验收产物。论文
分析将编号论文、赛题原文和专家评述分为三类证据：编号论文用于结构和文风统计，
赛题原文用于核对任务，专家评述只用于补充题型背景。全量页面采用文本层提取、
RapidOCR 和 Tesseract 复识相结合的流程，并对公式、表格和小字号图注回看原始栅格。
在此基础上形成十四类章节写作规范、A/B/C 题微观行文差异、创新策略、失分风险与
四周量化训练安排，同时记录双模式论文模板和写作审计 Skill 的实际完成路径。

\textbf{关键词：}数学建模竞赛；机理建模；组合优化；数据决策；论文写作；证据台账
\end{abstract}

\tableofcontents
\newpage

